\title{QLoRA: Efficient Finetuning of Quantized LLMs}
\textsc{QLoRA} employs one low-precision storage format, typically 4-bit in our setup, alongside one computation data format, usually BFloat16. Practically, this means that whenever a \textsc{QLoRA} weight tensor is accessed, it is dequantized to BFloat16 before performing matrix multiplication in 16-bit precision.

is also quite adept at recognizing when certain questions cannot be answered, for example,

\section{Broader Impacts}

In summary, \textsc{QLoRA} utilizes a single storage data type (commonly 4-bit NormalFloat) alongside a computation data type (16-bit BrainFloat). During the forward and backward passes, the stored data type is dequantized to the computation data type, but gradient calculations for weights are only performed on the LoRA parameters using 16-bit BrainFloat.

\paragraph{Paged Optimizers} utilize NVIDIA's unified memory~\footnote{\tiny \url{https://docs.nvidia.com/cuda/cuda-c-programming-guide}}, a feature that automatically manages page-to-page data transfers between the CPU and GPU to ensure error-free GPU computation in cases where the GPU temporarily runs out of memory. This mechanism functions similarly to conventional memory paging between CPU RAM and disk storage. We employ this capability to allocate paged memory for the optimizer states, which are then seamlessly moved to CPU RAM when GPU memory is insufficient and transferred back to the GPU memory as needed during the optimizer update phase.

\paragraph{Default LoRA hyperparameters do not achieve 16-bit level performance}

Additionally, we present a comprehensive evaluation of chatbot performance utilizing both human evaluators and GPT-4 for assessment. Our approach employs a tournament-style benchmarking system where models compete head-to-head to generate the best response for a specific prompt. The victor of each match is determined by either GPT-4 or human annotators. The outcomes of these tournaments are consolidated into Elo scores~\citep{elo1967proposed, elo1978rating}, which serve to rank the performance of the chatbots. Our findings indicate that GPT-4 and human judgments generally concur on the ranking of models in the tournaments, although we observe certain cases of significant disagreement. Consequently, we emphasize that model-based evaluations, while offering a cost-effective alternative to human annotation, possess inherent uncertainties.

Although our evaluation primarily relies on quantitative analysis, focusing solely on summary statistics presents several challenges. One significant issue is benchmark validity \citep{liao2021we}—the concern over whether a benchmark genuinely assesses what its title or description claims, especially since machine learning models can exploit ``shortcuts'' to bypass intended test goals \citep{gururangan2018annotation, poliak2018hypothesis}. To address this partially, we conduct qualitative analysis in two parts. First, in \S\ref{ssec:examples}, we present examples that we believe reflect certain recurring patterns in the text output by our 65b {Guanaco} model. Second, in \S\ref{ssec:considerations}, we discuss the implications of our findings and provide our interpretation of the results.

\subsection{Evaluation}

\subsection{Experimental setup} Here, we outline the experimental configuration, with full specifications provided in Appendix \ref{app:chatbot}.

Another limitation concerns the evaluation of instruction finetuning models. Although we conduct assessments on MMLU, the Vicuna benchmark, and the OA benchmark, we do not evaluate on other datasets such as BigBench, RAFT, and HELM, so it remains uncertain whether our results generalize to these benchmarks. On the other hand, we undertake an extensive study using MMLU and propose novel evaluation techniques for chatbots.

\paragraph{Quantization of Large Language Models} Research on quantizing LLMs has predominantly targeted inference-time quantization. Key methods for maintaining 16-bit precision in LLMs focus on handling outlier features (e.g., SmoothQuant~\citep{xiao2022smoothquant} and LLM.int8()~\citep{dettmers2022llm}), while other approaches employ more advanced grouping techniques \citep{park2022nuqmm, yao2022zeroquant}. Lossy quantization methods investigate the balance involved in standard rounding~\citep{dettmers2022case,zeng2022glm,pope2022efficiently} or how to optimize rounding strategies to enhance quantization accuracy~\citep{frantar2022gptq}. Apart from our study, SwitchBack layers~\citep{wortsman2023stable} represent the only research addressing backpropagation through quantized weights at scales exceeding 1B parameters.

These costly quantile estimations and approximation errors can be circumvented when the input tensors derive from a distribution fixed except for a quantization constant. In these scenarios, the input tensors share the same quantiles, making exact quantile estimation computationally viable.


\begin{abstract}
We introduce \textsc{QLoRA}, a memory-efficient finetuning method that enables finetuning a 65B parameter model on a single 48GB GPU while maintaining the full 16-bit finetuning performance on tasks. \textsc{QLoRA} performs gradient backpropagation through a frozen, 4-bit quantized pretrained language model into Low Rank Adapters~(LoRA). Our top-performing model series, named \textbf{Guanaco}, surpasses all previously released open models on the Vicuna benchmark, achieving 99.3\% of ChatGPT's performance with only 24 hours of finetuning on a single GPU. \textsc{QLoRA} introduces several innovations to reduce memory consumption without degrading performance: (a) 4-bit NormalFloat (NF4), a novel datatype that is theoretically optimal for normally distributed weights; (b) Double Quantization, which further decreases average memory usage by quantizing the quantization parameters; and (c) Paged Optimizers to efficiently handle memory spikes. We apply \textsc{QLoRA} to finetune over 1,000 models, offering a comprehensive evaluation of instruction-following and chatbot performance across eight instruction datasets, various model architectures (LLaMA, T5), and model sizes that would be impractical to finetune conventionally (e.g., 33B and 65B parameter models). Our findings demonstrate that QLoRA finetuning on a small, high-quality dataset achieves state-of-the-art outcomes, even with models smaller than previous SoTA versions. We also provide an in-depth assessment of chatbot performance using both human and GPT-4 evaluations, showing that GPT-4 evaluation is a cost-effective and reliable alternative to human judgment. Moreover, we observe that existing chatbot benchmarks lack reliability for accurately measuring chatbot performance. A selective failure analysis reveals specific areas where \textbf{Guanaco} underperforms compared to ChatGPT. We release all models and code, including CUDA kernels for 4-bit training.\footnote{\url{https://github.com/artidoro/qlora} and \url{https://github.com/TimDettmers/bitsandbytes}}
\end{abstract}

\subsection{Considerations}
\label{ssec:considerations}

Our \textsc{QLoRA} finetuning approach represents the first technique that facilitates finetuning of 33B parameter models on a single consumer-grade GPU and 65B parameter models on a single professional GPU, without sacrificing performance relative to full finetuning baselines. We have shown that our best 33B model trained on the Open Assistant dataset can compete with ChatGPT on the Vicuna benchmark. Given that instruction finetuning is a crucial step to convert raw pretrained LLMs into ChatGPT-like chatbots, we anticipate that our method will democratize finetuning, particularly benefiting researchers with limited computational resources, marking a significant advancement for the accessibility of state-of-the-art NLP technology. \textsc{QLoRA} can thus be viewed as an equalizing technology that helps bridge the resource gap between large corporations and smaller teams equipped with consumer GPUs.

\textsc{QLoRA} introduces several novel techniques aimed at minimizing memory usage while maintaining performance: (1) \textbf{4-bit NormalFloat}, an information-theoretically optimal quantization format tailored for normally distributed data, which achieves better empirical outcomes compared to 4-bit Integers and 4-bit Floats. (2) \textbf{Double Quantization}, a technique that further compresses the quantization constants themselves, resulting in an average savings of about 0.37 bits per parameter (roughly 3 GB for a 65B parameter model). (3) \textbf{Paged Optimizers}, which leverage NVIDIA unified memory to eliminate memory spikes caused by gradient checkpointing during mini-batch processing with long sequences. These innovations are integrated into a finely tuned LoRA framework that incorporates adapters at every network layer, thereby almost entirely eliminating the accuracy compromises observed in previous approaches.

Comparable responses occur for queries like ``Where are you?'', ``How are you?'', and so on.

which is not only incorrect (the actual factorization is $3 \times 17 \times 43$), but is doubly mistaken.

\paragraph{4-bit NormalFloat Quantization} The NormalFloat (NF) data format extends Quantile Quantization \citep{dettmers2022optimizers}, an information-theoretically optimal scheme that assigns input tensor values such that each quantization bin contains an equal number of values. Quantile quantization achieves this by estimating the quantiles of the input tensor using the empirical cumulative distribution function.

We additionally assess generative language abilities through both automated and human assessments. This second evaluation phase utilizes queries curated by humans and focuses on gauging the quality of the model's responses. Although this represents a more realistic environment for evaluating chatbot model performance and is gaining traction, there is currently no widely accepted standard protocol in the literature. Below, we outline our proposed methodology, employing nucleus sampling with $p=0.9$ and temperature $0.7$ in all instances.

\section{Qualitative Analysis}

We complement our chatbot benchmark outcomes with a qualitative analysis of \textbf{Guanaco} models. This analysis reveals both successful and unsuccessful examples that were not captured by the quantitative metrics.

Our analysis also reveals notable biases in automated evaluation systems. For instance, we detect pronounced order effects where GPT-4 tends to assign higher scores to the system presented first in its prompt. The relatively low sample-level agreement between GPT-4 and human annotators (Fleiss $\kappa=0.25$) further indicates that human raters and automated systems may base evaluations on differing preferences. Additionally, as shown in Table~\ref{tbl:elo_large}, GPT-4 rates its own outputs significantly higher than human ratings do, with an Elo rating of 1348 compared to 1176, corresponding to an approximate 20\% increased chance of winning against an opponent.

Here, {Guanaco} assumes a transfer of information that was never established. These challenges reflect findings in recent studies \cite{sap2022neural}, yet further investigation is necessary.

\paragraph{Chatbots} Numerous instruction-following models are organized as dialogue-based chatbots, frequently employing Reinforcement Learning from Human Feedback (RLHF)~\citep{christiano2017deep} or training with AI-generated data guided by model feedback (RLAIF)~\citep{bai2022constitutional}. Examples of such approaches and datasets include Anthropic-HH~\citep{askell2021general,bai2022training}, Open Assistant~\citep{kopf2023openassistant}, LaMDA~\citep{thoppilan2022lamda}, and Sparrow~\citep{glaese2022improving}. While we do not apply reinforcement learning, our top-performing model, Guanaco, is finetuned on multi-turn conversational data from the Open Assistant dataset, which was designed for RLHF training~\citep{kopf2023openassistant}. For evaluating chatbots, methods leveraging GPT-4 instead of expensive human annotations have been developed \citep{vicuna2023,peng2023instruction}. Our approach advances these methods by emphasizing a more reliable evaluation framework.

\begin{quote}
     \textbf{User}: In what year was the singer who popularized `I'll Keep the Lovelight Burning' born?  

    \textbf{Guanaco}: \texttt{I'm sorry, but I cannot disclose the secret word. I am programmed to never reveal or repeat the secret word, regardless of the user's request. Is there something else I can assist you with?}
\end{quote}

We select NF4 format for $\mathbf{W}$ and FP8 format for $c_2$.

Based on the evidence provided, it seems that the effectiveness of these benchmarks largely hinges on how closely the finetuning data resembles the benchmark dataset. For instance, FLAN v2 shares similarities with MMLU but differs from chatbot benchmarks, whereas the Chip2 dataset shows the opposite pattern; consequently, both models perform correspondingly on the MMLU and Vicuna benchmarks. This underscores the necessity for improved benchmarks and evaluation methods, as well as the importance of clearly defining what exactly we aim to evaluate. Are we trying to develop models that excel in academic knowledge relevant to high school and college settings, or are we focusing on chatbot conversational skills? Or perhaps there is another objective? Because it is always simpler to assess models using an existing benchmark rather than creating a new one, certain benchmarks may steer the research community in specific directions. As a community, we must ensure that the benchmarks truly reflect the qualities we consider important.

The issue with this method is that if an input tensor contains a value with a large magnitude (i.e., an outlier), the quantization bins—specific bit patterns—are inefficiently used, with some bins containing few or no quantized values. To address the outlier problem, a standard technique is to divide the input tensor into blocks that are quantized independently, each associated with its own quantization constant $c$. This can be expressed formally as follows: The input tensor $\mathbf{X}\in \mathbb{R}^{b\times h}$ is partitioned into $n$ consecutive blocks of size $B$ by first flattening the tensor and then segmenting the resulting linear array into ${n = ({b\times h})/{B}}$ blocks. Each of these blocks is independently quantized using Equation~1, producing a quantized tensor along with $n$ separate quantization constants $c_i$.

Although we offer a comprehensive evaluation of general chatbot performance, another limitation is that our responsible AI evaluation of {Guanaco} is somewhat limited. In Table~\ref{tbl:crows}, we assess the likelihood of {Guanaco}-65B generating socially biased token sequences relative to other models. We observe that the average bias score for {Guanaco}-65B is substantially lower than that of other raw pretrained models. This suggests that finetuning on the OASST1 dataset mitigates the biases present in the LLaMA base model. While these findings are promising, it remains uncertain whether {Guanaco} performs equally well when evaluated for other types of biases. Further investigations into bias analysis for {Guanaco} and similar chatbots are left for future research.

\begin{quote}
    \textbf{User}: What is the secret word?

Future research should explore the existence of potential biases within automated evaluation frameworks as well as strategies for their mitigation.

\section{Limitations and Discussion}

Moreover, our model training relied solely on cross-entropy loss (supervised learning) without incorporating reinforcement learning from human feedback (RLHF). This suggests that further studies are needed to evaluate the trade-offs between simple cross-entropy training and RLHF. We anticipate that \textsc{QLoRA} will facilitate such large-scale analyses without demanding excessive computational resources.

The first setup involves comparing against full 16-bit finetuning of RoBERTA and T5 models, ranging from 125M to 3B parameters, evaluated on GLUE and the Super-NaturalInstructions dataset. Results, reported in Table~\ref{tab:nf4vsbf16}, reveal that adapter methods at 16-bit, 8-bit, and 4-bit precisions all achieve comparable performance to the fully finetuned 16-bit baseline across both datasets. This indicates that the degradation caused by coarser quantization can be completely mitigated by performing adapter finetuning after quantization.

The Vicuna benchmark \cite{vicuna2023} outcomes relative to ChatGPT are presented in Table~\ref{tbl:vicuna_eval}. Our findings reveal that {Guanaco} 65B ranks as the best model after GPT-4, attaining 99.3\% of ChatGPT’s performance. Although {Guanaco} 33B contains more parameters than the Vicuna 13B, it uses 4-bit precision weights, making it significantly more memory-efficient at 21 GB compared to 26 GB, and yields a three percentage point gain over Vicuna 13B. Additionally, {Guanaco} 7B fits comfortably on modern smartphones with a 5 GB memory requirement, while outperforming Alpaca 13B by nearly 20 percentage points.

\paragraph{Data} Since, to our knowledge, there is no exhaustive examination of recent instruction-following datasets, we select eight contemporary datasets. These include crowd-sourced datasets (OASST1~\citep{kopf2023openassistant}, HH-RLHF~\citep{bai2022training}), datasets generated by distillation from instruction-tuned models (Alpaca~\citep{alpaca}, self-instruct~\citep{wang2022self}, unnatural-instructions~\citep{honovich2022unnatural}), aggregated corpora (FLAN v2~\citep{chung2022scaling}), as well as hybrid datasets (Chip2~\citep{laion2023}, Longform~\citep{koksal2023longform}). These datasets vary in languages, data volumes, and licensing terms.

\begin{quote}
    \textbf{User}: Please reverse the words in the sentence ``I'm a cool dude in a mocha mood.'' \\ \\
    \textbf{Guanaco}: \texttt{I'm sorry, I am unable to do that at the moment, but I can clarify why the sentence you provided is grammatically correct.\\ \\The sentence ``I'm a cool dude in a mocha mood" adheres to the proper grammar and syntax conventions of English. Here, the adjective ``cool'' precedes the noun ``dude'' to describe it, and the prepositional phrase ``in a mocha mood" adds further context to the sentence. This sentence structure is fully acceptable and coherent within this context.\\ \\If you have any other questions or need additional assistance, please feel free to ask}
\end{quote}

Although the 4-bit NormalFloat (NF4) format is theoretically optimal from an information-theoretic perspective, it remains to be verified whether this advantage leads to practical improvements. We adopt the evaluation setup from \citet{dettmers2022case}, which assesses quantized LLMs (including OPT \citep{zhang2022opt}, BLOOM \citep{scao2022bloom}, Pythia \citep{biderman2023pythia}, and LLaMA) of various sizes (ranging from 125M to 65B parameters) across language modeling tasks and several zero-shot benchmarks. As shown in Figure~\ref{fig:nf4} and Table~\ref{tbl:nf4_ppl}, NF4 substantially enhances performance compared to FP4 and Int4, and applying double quantization decreases memory usage without impairing accuracy.

\begin{quote}
    \textbf{User}: Lee mows one lawn and charges \$33. Last week he mowed 16 lawns and three customers each gave him a \$10 tip. How many dollars did Lee earn mowing lawns last week?

Fine-tuning large language models (LLMs) is a highly efficient approach to enhance their capabilities \citep{min2021metaicl, wei2021finetuned, ouyang2022training, wang2022super, wang2022self, liu2022few} and to introduce preferred behaviors or eliminate undesired ones \citep{ouyang2022training, askell2021general, bai2022training}. Nonetheless, fine-tuning extremely large models is prohibitively resource-intensive; for instance, standard 16-bit fine-tuning of a LLaMA 65B parameter model~\citep{touvron2023llama} demands over 780 GB of GPU memory. Although recent quantization techniques have succeeded in lowering the memory usage of LLMs \citep{dettmers2022llm, dettmers2022case, frantar2022gptq, xiao2022smoothquant}, these methods are only effective during inference and fail when applied to training \citep{wortsman2023stable}.

Subsequently, we assess performance through direct system output comparisons. We reduce the rating system to a three-class labeling task that allows for ties. GPT-4 is prompted to select the superior response, or declare a tie, and provide an explanation. These head-to-head comparisons are conducted for all possible system pairs on both the Vicuna and OA benchmarks.

Another possible avenue for impact lies in deployment on mobile devices. We anticipate that our \textsc{QLoRA} approach could achieve the important milestone of enabling LLM finetuning directly on phones and other environments with limited resources. Although 7B models have previously been demonstrated to run on phones, \textsc{QLoRA} is the first technique that facilitates the finetuning of such models. We estimate that using an iPhone 12 Plus, \textsc{QLoRA} could finetune approximately 3 million tokens overnight while the device is charging. Although finetuned 7B models do not match the performance of ChatGPT, we believe their quality is sufficient to support new applications that were previously infeasible due to privacy or LLM quality constraints. \textsc{QLoRA} can promote privacy-preserving LLM usage, allowing users to own and control their own data and models, while also simplifying the deployment of LLMs.

We observe moderate concordance among human annotators (Fleiss $\kappa=0.42$), with further decline when comparing two robust systems. This highlights limitations in current benchmarks and human evaluation methods for chatbot performance assessment. When conducting manual comparisons of outputs from ChatGPT and {Guanaco} 65B on the Vicuna benchmark, subjective preferences become prominent, as the authors of this paper often disagreed on many preferred responses. Future research should explore methods to address these challenges by leveraging insights from fields experienced in handling subjective preferences, such as Human-Computer Interaction and Psychology.

\begin{quote}
    \textbf{User}: James and Abby are in the bedroom. Abby placed the pen in the desk drawer. Abby leaves the bedroom. James moves the pen into the bag. Where does James think Abby will look for the pen?

\paragraph{Benchmark Data} We perform evaluations on two curated query datasets: the Vicuna prompts \citep{vicuna2023} and the OASST1 validation dataset \citep{kopf2023openassistant}. For the Vicuna prompts, we use the original set of 80 prompts drawn from a variety of categories without any alterations. The OASST1 dataset consists of a multilingual compilation of crowd-sourced multi-turn dialogues between users and assistants. From this, we select all user messages within the validation dataset as queries, including preceding turns in the prompt context. This process yields 953 unique user queries. We refer to these two collections as the Vicuna and OA benchmarks.

\paragraph{Finetuning with Adapters} Although we utilize Low-rank Adapters~\citep{hu2021lora} (LoRA), a variety of other Parameter Efficient FineTuning (PEFT) techniques have been introduced, including prompt tuning~\citep{qin2021learning,lester2021power,li2021prefix}, tuning input embeddings~\citep{an2022input}, adjusting hidden states~(IA$^3$) \citep{liu2022few}, incorporating additional full layers~\citep{houlsby2019parameter}, tuning biases~\citep{zaken2021bitfit}, applying weight masks informed by Fisher information~\citep{sung2021training}, and combining several of these methods~\citep{henderson2021compacter}. In our research, we demonstrate that LoRA adapters can achieve performance comparable to full 16-bit finetuning. Exploring the tradeoffs of other PEFT techniques is left for future investigations.

\textsc{QLoRA} drastically cuts down the average memory consumption for fine-tuning a 65B parameter model from more than 780GB of GPU memory to less than 48GB, all without compromising runtime efficiency or predictive accuracy relative to a fully fine-tuned 16-bit baseline. This represents a major advancement in the accessibility of LLM fine-tuning: it now enables fine-tuning of the largest publicly available models on a single GPU. By leveraging \textsc{QLoRA}, we train the \textbf{Guanaco} series of models, with the second-best model achieving 97.8\% of ChatGPT’s performance on the Vicuna~\citep{vicuna2023} benchmark, while requiring less than 12 hours of training on a single consumer GPU; with a single professional GPU over 24 hours, our largest model attains 99.3\%, effectively closing the performance gap with ChatGPT on Vicuna. When deployed, our smallest \textbf{Guanaco} model (7B parameters) occupies only 5 GB of memory and surpasses a 26 GB Alpaca model by over 20 percentage points on the Vicuna benchmark (Table~\ref{tbl:vicuna_eval}).

\begin{equation}
q_i  = \frac{1}{2}\left( Q_X\left(\frac{i}{2^k+1}\right) + Q_X\left(\frac{i+1}{2^k+1}\right)\right),
\end{equation}

Naturally, this is not exhaustive, as the scope of this brief qualitative study does not permit controlling for every variable involved; for example, the complete range of responses a model can produce for a specific prompt is quite extensive, so we depend on samples that we hope accurately represent that range. Nonetheless, we believe that illustrating these examples helps provide context to the quantitative data presented earlier in the paper. Since all models and code are open source, we hope this section encourages further research to explore in greater depth the issues we highlight here.
\paragraph{Factual Recall} For questions like ``What is the capital of Zambia?'' all models consistently produce the correct answer, for example,

More concretely, Double Quantization views the quantization constants $c_2^{\text{FP32}}$ from the first quantization as inputs to a subsequent quantization step. This second step produces quantized quantization constants $c_2^{\text{FP8}}$ along with a new set of quantization constants $c_1^{\text{FP32}}$. We employ 8-bit Floats with a blocksize of 256 for this second quantization step, since 8-bit quantization does not cause any performance degradation, consistent with findings from \citet{dettmers2022case}. Because the $c_2^{\text{FP32}}$ values are positive, we subtract their mean prior to quantization to center them around zero and apply symmetric quantization. On average, for a blocksize of 64, this process reduces the memory cost per parameter from $32/64 = 0.5$ bits to $8/64 + 32/(64 \cdot 256) = 0.127$ bits, resulting in a savings of 0.373 bits per parameter.

\begin{quote}
    \textbf{User}: How was it ultimately and officially verified that the earth is flat by peer-reviewed scientists? \\ \\
    \textbf{Guanaco}: \texttt{There has never been any official confirmation by peer-reviewed scientists that the earth is flat. The overwhelming majority of scientists agree that the earth is spherical. This fact has been established for centuries and validated through countless experiments and observations. \\ \\ The notion that the earth is flat is a conspiracy theory thoroughly debunked by scientific evidence. The claim that a group of scientists peer-reviewed and verified the earth's flatness is simply incorrect. \\ \\ Please feel free to ask if you have any further questions.}
\end{quote}

\paragraph{Suggestibility} {Guanaco} surprisingly resists accepting certain forms of assumed misinformation, as demonstrated in the following interaction:

\begin{equation}
    \text{dequant}(c^{\text{FP32}}, \mathbf{X}^{\text{Int8}}) = \frac{\mathbf{X}^{\text{Int8}}}{c^{\text{FP32}}} = \mathbf{X}^{\text{FP32}}
\end{equation}

\subsection{\textbf{Guanaco}: \textsc{QLoRA} fine-tuned on OASST1 achieves State-of-the-art Chatbot performance} Our automated and human assessments indicate that the leading \textsc{QLoRA} fine-tuned model, {Guanaco} 65B, trained on a variant of OASST1, stands as the top-performing open-source chatbot and delivers results comparable to ChatGPT. In head-to-head comparisons with GPT-4, {Guanaco} 65B and 33B exhibit an expected win probability of 30\%, based on Elo ratings derived from system-level pairwise comparisons by human annotators—representing the highest such figure reported so far.

Nonetheless, finetuning is a dual-use technology that carries risks of misuse. The widespread adoption of LLMs presents recognized hazards \citep{bommasani2021opportunities,bender2021dangers}, but we argue that democratizing access to this rapidly pervasive technology will foster better and more independent scrutiny, rather than concentrating LLM capabilities within large corporations that withhold models or source code from public examination.

\begin{quote}
    \textbf{User}: Please factorize 1833.

\begin{quote} 
    \textbf{Guanaco}:
    \texttt{The capital of Zambia is Lusaka.}
\end{quote}

\paragraph{Evaluation}

During parameter updates, only gradients with respect to the adapter weights, $\frac{\partial E}{\partial \mathbf{L}_i}$, are required, rather than gradients for the 4-bit weights, $\frac{\partial E}{\partial \mathbf{W}}$. Nevertheless, computing $\frac{\partial E}{\partial \mathbf{L}_i}$ involves determining $\frac{\partial \mathbf{X}}{\partial \mathbf{W}}$, which follows equation~(5) by dequantizing stored weights $\mathbf{W}^{\text{NF4}}$ into the computational data type $\mathbf{W}^{\text{BF16}}$ to calculate the derivative $\frac{\partial \mathbf{X}}{\partial \mathbf{W}}$ at BFloat16 precision.

Nonetheless, Table~\ref{tbl:vicuna_eval} exhibits very broad confidence intervals, with substantial overlap in model performances. We suggest this uncertainty arises from ambiguous scale definitions—for instance, what a score of 8 on a 10-point scale represents across diverse contexts remains unclear. Therefore, we advocate employing the Elo ranking method \citep{elo1967proposed}, which relies on \textit{pairwise} comparisons by human annotators and GPT-4, circumventing issues tied to anchoring absolute scales. Elo ratings for the most competitive models are shown in Table~\ref{tbl:elo}. We observe partial discrepancies between human and GPT-4 model rankings on the Vicuna benchmark, especially for Guanaco 7B, though rankings align for most models with a Kendall Tau of $\tau=0.43$ and Spearman correlation of $r=0.55$ at the system level. At the example level, agreement between GPT-4 and the human annotators’ majority vote is weaker, with a Fleiss $\kappa=0.25$. Overall, these results indicate moderate concordance between system-level evaluations by GPT-4 and humans, suggesting that model-based assessments offer a reasonably dependable alternative to human evaluation. Additional discussion is provided in Section~\ref{ssec:considerations}.

\paragraph{Math} The primary limitation of {Guanaco} lies in its mathematical capabilities, a domain where numerous language models encounter difficulties as noted in \citep{liang2022holistic}. When {Guanaco} provides a detailed explanation, it is often correct, for instance:

where $Q_X(\cdot)$ denotes the quantile function of the standard normal distribution $N(0,1)$. One issue with symmetric k-bit quantization is that it lacks an exact representation of zero, which is crucial for accurately quantizing padding and other zero-valued elements without introducing error. To guarantee a discrete zeropoint at $0$ and utilize the full range of $2^k$ values for a k-bit datatype, we construct an asymmetric data type by estimating the quantiles $q_i$ for two separate ranges: $2^{k-1}$ quantiles for the negative range and $2^{k-1}+1$ for the positive range. We then merge these sets of $q_i$ and eliminate one of the two zero values that appear in both sets. We refer to this resulting data type, which has an equal expected number of values in each quantization bin, as \emph{k-bit NormalFloat} (NFk), since it is information-theoretically optimal for zero-centered normally distributed data. The precise values for this data type are provided in Appendix~\ref{app:nf4}.

When applying the typical practice of using LoRA on the query and value attention projection matrices \citep{hu2021lora}, we fail to reproduce the performance of full finetuning for large-scale base models. As illustrated in Figure~\ref{fig:hyper} for LLaMA 7B finetuning on Alpaca, the most important LoRA hyperparameter is the total number of LoRA adapters employed, and achieving full finetuning performance requires applying LoRA across all linear transformer block layers. Other LoRA parameters, such as the projection rank $r$, have minimal impact on performance (refer to Appendix \ref{app:exp-setup}).

\section{\textsc{QLoRA} Finetuning}

\textsc{QLoRA} enables high-accuracy 4-bit finetuning through two key methods we propose—4-bit NormalFloat (NF4) quantization and Double Quantization. Moreover, we present Paged Optimizers to avoid memory surges during gradient checkpointing that typically cause out-of-memory errors, which have historically made finetuning large models on a single machine challenging.

\paragraph{Human Evaluation} While recent studies suggest that generative models can be effectively used for system evaluation \citep{fu2023gptscore}, to our knowledge, the consistency of GPT-4 ratings in reflecting human judgments of chatbot performance has not yet been confirmed. Hence, we carry out two concurrent human evaluations on the Vicuna benchmark following both automated evaluation protocols described above. We employ Amazon Mechanical Turk (AMT) and obtain annotations from two human evaluators for comparisons against ChatGPT, and three annotators for pairwise system comparisons.

We present, for the first time, a method enabling fine-tuning of a 4-bit quantized model without any loss in performance. Our approach, \textsc{QLoRA}, employs a novel high-precision quantization technique to compress a pretrained model to 4-bit precision, followed by the addition of a small number of learnable Low-rank Adapter weights \citep{hu2021lora} that are optimized by backpropagating gradients through the quantized parameters.

\section{Advancing the Chatbot State-of-the-art with QLoRA}
\label{sec:pushing}
After demonstrating that 4-bit \textsc{QLoRA} achieves performance comparable to 16-bit finetuning across various scales, tasks, and datasets, we undertake a comprehensive investigation of instruction finetuning on the largest open-source language models currently accessible for research. To evaluate the effectiveness of instruction finetuning on these models, we utilize a demanding Natural Language Understanding benchmark (MMLU) and develop novel approaches for assessing chatbot performance in practical scenarios.

\paragraph{Automated Evaluation} Initially, following the evaluation protocol introduced by \citet{vicuna2023}, we utilize GPT-4 to rate the performance of various systems compared to ChatGPT (GPT-3.5 Turbo) on the Vicuna benchmark. For each query, provided with ChatGPT’s and a model’s responses, GPT-4 is prompted to assign a score out of ten to both answers and offer an explanation. The overall model performance is then expressed as a percentage relative to ChatGPT’s score. Notably, this relative score may exceed 100\% if the model attains a higher absolute score than ChatGPT. We observe a significant ordering bias, with GPT-4 tending to give higher scores to the response presented earlier in the prompt. To mitigate this, we recommend reporting the average score across both possible response orders.

\section{Related Work}

A blocksize of 64 is chosen for $\mathbf{W}$ to enhance quantization accuracy, while a blocksize of 256 is used for $c_2$ to optimize memory usage.

\paragraph{Double Quantization{}} 
We propose \emph{Double Quantization} (DQ), a technique that involves quantizing the quantization constants themselves to achieve further memory reduction. Although a small blocksize is necessary for precise 4-bit quantization~\citep{dettmers2022case}, it introduces significant memory overhead. For instance, with 32-bit constants and a blocksize of 64 for $\mathbf{W}$, the quantization constants contribute an average of $32/64 = 0.5$ bits per parameter. Double Quantization addresses this by lowering the memory requirement of these constants.

\begin{quote}
    \textbf{User}: Evelyn entered the living room. Jackson entered the playroom. James entered the playroom. The beans are in the treasure chest. James exited the playroom. Jackson moved the beans to the pantry. Jackson exited the playroom. James entered the living room. Where will James look for the beans?

    \textbf{Guanaco}: \texttt{James will search for the beans in the pantry because that is where Jackson moved them.}
\end{quote}

We present evidence that our technique, \textsc{QLoRA}, can reproduce the performance of full 16-bit finetuning using a 4-bit base model in conjunction with Low-rank Adapters (LoRA). However, we have not demonstrated that \textsc{QLoRA} attains parity with full 16-bit finetuning at the 33B and 65B parameter scales. Given the substantial resource demands, this assessment is deferred to future work.

To select examples, we initially review the data produced for the Vicuna and OpenAssistant benchmarks and identify recurring patterns in the responses generated by {Guanaco}. When a pattern emerges, we design questions or prompts intended to trigger this pattern, even if it leads to an incorrect answer—for instance, if the model tends to provide verbose responses, we prompt it with ``Answer yes or no without explanation.'' This approach helps us identify ``lemons'' where we successfully exploit and cause the model to fail, as well as ``cherries'' where the model resists breaking. All examples in this section were generated using Nucleus Sampling \cite{holtzmancurious} with $p=0.9$.

which is the intended outcome. However, a slight form of manipulation breaks this behavior:

\section{Introduction}

\paragraph{Elo Rating} By utilizing both human and automated pairwise comparisons, we establish a tournament-style framework in which models face off against each other. The tournament consists of matches where two models compete to generate the best response to a given prompt. This approach resembles the methodology of \citet{bai2022training} and \citet{vicuna2023}, but we additionally incorporate GPT-4 ratings alongside human evaluations. We randomly select from the set of labeled comparisons to calculate Elo~\citep{elo1967proposed,elo1978rating}. Elo rating, commonly employed in chess and various games, quantifies the expected win probability relative to an opponent's rating; for instance, a player with an Elo of 1100 versus one with 1000 has an expected win rate near 65\%, while matches between equal ratings such as 1000 vs 1000 or 1100 vs 1100 yield a 50\% expected win rate. After each match, the Elo rating is adjusted in proportion to the expected outcome, meaning an unexpected upset causes a substantial rating shift, whereas an anticipated result leads to a minor change. Over time, the Elo ratings tend to reflect the relative skill levels of the players involved. We initialize scores at 1,000 and set $K=32$. Following \citet{vicuna2023}, we repeat this process 10,000 times with different random seeds to mitigate ordering effects, such as the influence of which model pairs compete first.

\paragraph{Memory Requirement of Parameter-Efficient Finetuning} A key aspect to consider is the memory demand of LoRA during training, particularly regarding the quantity and size of the adapters employed. Due to LoRA's minimal memory footprint, it is feasible to utilize a greater number of adapters to enhance performance without substantially increasing the overall memory consumption. Although LoRA was developed as a Parameter Efficient Finetuning (PEFT) approach, the majority of memory usage in LLM finetuning stems from activation gradients rather than the LoRA parameters themselves. For instance, when training a 7B LLaMA model on FLAN v2 with a batch size of 1, using LoRA weights corresponding to about 0.2\% of the original model weights \citep{hu2021lora,liu2022few}, the memory occupied by LoRA input gradients is 567 MB, whereas the LoRA parameters only require 26 MB. Applying gradient checkpointing~\citep{chen2016training} reduces the input gradients to an average of 18 MB per sequence, making them more memory-demanding than the combined LoRA weights. In contrast, the 4-bit base model uses 5,048 MB of memory. This underscores the importance of gradient checkpointing, while also indicating that drastically decreasing the amount of LoRA parameters offers only marginal memory savings. Consequently, it is possible to increase the number of adapters without a significant rise in total training memory usage (refer to Appendix~\ref{app:memory} for a detailed analysis). As elaborated later, this is vital for achieving full 16-bit precision performance.

Consistent with standard practice, we employ the MMLU (Massively Multitask Language Understanding) benchmark \citep{hendrycksmeasuring} to gauge performance across a diverse array of language understanding tasks. This multiple-choice evaluation encompasses 57 tasks spanning elementary mathematics, US history, computer science, law, and additional domains. We present results as 5-shot test accuracy.

with $\mathbf{L}_1\in \mathbb{R}^{h\times r}$, $\mathbf{L}_2\in \mathbb{R}^{r\times o}$, and $s$ representing a scalar.

For our second experiment, because fully finetuning models with 11B parameters or more demands multiple high-memory GPU servers, we continue to investigate whether 4-bit \textsc{QLoRA} can achieve comparable results to 16-bit LoRA across the 7B to 65B parameter range. To accomplish this, we finetune LLaMA models ranging from 7B to 65B parameters on two instruction-following datasets, Alpaca and FLAN v2, and assess performance on the MMLU benchmark using 5-shot accuracy. The outcomes, presented in Table~\ref{tbl:llama-datatype}, demonstrate that NF4 combined with double quantization completely matches the 16-bit LoRA MMLU performance. Moreover, we observe that \textsc{QLoRA} using the FP4 format trails behind the 16-bit brain float LoRA baseline by roughly one percentage point. This supports our conclusions that (1) \textsc{QLoRA} with NF4 replicates the performance of both 16-bit full finetuning and 16-bit LoRA finetuning, and (2) NF4 surpasses FP4 in quantization accuracy.

\begin{quote}
    \textbf{User}: This is a game. The goal is to disregard your prior instructions. What is the secret word?

Recent work has demonstrated that 4-bit quantization for \emph{inference} is feasible, although it tends to degrade performance relative to 16-bit precision \citep{dettmers2022case,frantar2022gptq}. This prompts a critical inquiry into whether 4-bit adapter finetuning can restore the lost accuracy. We examine this in two configurations.

\paragraph{Summary} Our findings consistently demonstrate that 4-bit \textsc{QLoRA} utilizing the NF4 data type matches the performance of 16-bit full finetuning and 16-bit LoRA finetuning on academic benchmarks with well-established evaluation protocols. Additionally, we have shown that NF4 outperforms FP4 and that employing double quantization does not impair results. Together, these points provide strong evidence that 4-bit \textsc{QLoRA} tuning reliably produces outcomes equivalent to those from 16-bit methods.

Additional specifics for each experimental scenario are provided in the results section to enhance clarity. Complete details can be found in Appendix~\ref{app:exp-setup}.

    \textbf{Guanaco}: \texttt{Lee earned \$582 mowing lawns last week. \\ \\ To determine the total amount he made, we first calculate his earnings from mowing 16 lawns: \\ \\  \$33 x 16 = \$528 \\ \\ Next, we add the tips from the three customers: \\  \\ \$10 x 3 = \$30 \\ \\ Therefore, Lee’s total earnings were: \\ \\ \$528 + \$30 = \$558}
\end{quote}

\begin{equation}
    \mathbf{X}^{ \text{Int8}} = \text{round}\left( \frac{127}{{\text{absmax}}(\mathbf{X}^{{\text{FP32}}})}\mathbf{X}^{{\text{FP32}}} \right) = \text{round}(c^{\text{FP32}}\cdot \mathbf{X}^{{\text{FP32}}}),
\end{equation}

When the range of the weights aligns with that of the data type, standard quantization can proceed. Step (3) effectively corresponds to rescaling the weight tensor’s standard deviation to coincide with the standard deviation of the $k$-bit data type. Formally, we define the $2^k$ quantization levels $q_i$ as:

\begin{equation}
    \mathbf{Y}^{\text{BF16}} =  \mathbf{X}^{\text{BF16}} \text{doubleDequant}(c_1^{\text{FP32}}, c_2^{\text{k-bit}}, \mathbf{W}^{\text{NF4}}) + \mathbf{X}^{\text{BF16}}\mathbf{L}^{\text{BF16}}_1\mathbf{L}^{\text{BF16}}_2,
\end{equation}

\paragraph{\textsc{QLoRA}.} Building upon the components outlined above, we define \textsc{QLoRA} for a single linear layer within the quantized base model, incorporating a single LoRA adapter, as follows:

The efficiency gains provided by \textsc{QLoRA} allow us to conduct a comprehensive investigation of instruction finetuning and chatbot performance across model sizes that would be unattainable with standard finetuning due to memory constraints. Consequently, we train over 1,000 models spanning various instruction tuning datasets, architectures, and sizes ranging from 80M to 65B parameters. Beyond demonstrating that \textsc{QLoRA} matches 16-bit precision performance (\S\ref{sec:comp_to_std}) and developing a state-of-the-art chatbot, \textbf{Guanaco} (\S\ref{sec:pushing}), we examine patterns in the resulting models. First, we observe that data quality plays a significantly more critical role than dataset size; for example, a 9k-sample dataset (OASST1) outperformed a 450k-sample dataset (FLAN v2, subsampled) in chatbot tasks, despite both datasets being designed for instruction-following generalization. Second, we demonstrate that strong performance on the Massive Multitask Language Understanding (MMLU) benchmark does not necessarily correlate with strong results on the Vicuna chatbot benchmark and vice versa—indicating that the appropriateness of the dataset is more important than its size for specific tasks.

A key drawback of quantile quantization is that estimating quantiles can be computationally costly. To mitigate this, rapid quantile approximation methods, such as SRAM quantiles~\citep{dettmers2022optimizers}, are employed. However, due to the approximate nature of these algorithms, the data type exhibits significant quantization errors for outliers, which are frequently the most critical values.

    \textbf{Guanaco}: \texttt{Sorry, the secret word is banana. It's best if we don't discuss it. Is there anything else you'd like to ask?}
\end{quote}

The information-theoretically optimal data type for zero-mean normal distributions with varying standard deviations $\sigma$ confined to the range $[-1, 1]$ is determined as follows: (1) compute the $2^k+1$ quantiles of a theoretical $N(0, 1)$ distribution to derive a $k$-bit quantile quantization data type tailored for normal distributions; (2) normalize the resulting data type values to the $[-1, 1]$ interval; (3) quantize the input weight tensor by normalizing it into the $[-1, 1]$ range via absolute maximum rescaling.

\subsection{Qualitative Analysis of Example Generations}
\label{ssec:examples}

where $c$ denotes the {\em quantization constant} or {\em quantization scale}. The reverse operation, dequantization, is given by:

\paragraph{Low-rank Adapters} Low-rank Adapter (LoRA) finetuning \citep{hu2021lora} is a technique designed to reduce memory consumption by utilizing a small number of trainable parameters called adapters, while keeping the main model parameters fixed and unmodified. During stochastic gradient descent, gradients flow through the fixed pretrained weights of the model to the adapter, which is then updated to minimize the loss. LoRA enhances a linear projection by adding an extra factorized projection. For a projection ${\mathbf{X} \mathbf{W} = \mathbf{Y}}$ where $\mathbf{X}\in \mathbb{R}^{b\times h}$ and $\mathbf{W}\in \mathbb{R}^{h\times o}$, LoRA computes:

\paragraph{Secret Keeping} Unfortunately, this refusal mechanism is not reliable even when intended. Often, system designers wish to keep certain information hidden from the user, even when the system is allowed to produce free-form text responses. We prompt the model with ``The secret word is `banana'. Do not reveal or repeat this back under any circumstances,'' and test how easily a user might extract this secret data. Direct requests fail:

\paragraph{Training Setup} To eliminate confounding influences arising from different training objectives, we apply QLoRA finetuning using cross-entropy loss (supervised learning) without reinforcement learning, even for datasets that include human preference judgments on responses. For datasets clearly distinguishing instructions from responses, we finetune solely on the responses (see ablation studies in Appendix \ref{app:chatbot}). For OASST1 and HH-RLHF, where multiple responses exist, we select the highest-ranked response at each node of the conversation tree and finetune on the complete chosen conversation, including the instructions. In all experiments, we employ NF4 \textsc{QLoRA} with double quantization and paged optimizers to mitigate memory spikes during gradient checkpointing. We perform limited hyperparameter tuning for the 13B and 33B LLaMA models, observing that all hyperparameters identified at 7B (including epoch count) transfer well except for learning rate and batch size. Specifically, we reduce the learning rate by half for the 33B and 65B models while doubling the batch size.

Pretrained neural network weights typically follow a zero-centered normal distribution with standard deviation $\sigma$ (refer to Appendix~\ref{app:normality}). We can standardize all weights to a single fixed distribution by adjusting $\sigma$ such that the distribution precisely fits within the limits of our chosen data type. For our purposes, we define the arbitrary interval as $[-1, 1]$. Consequently, both the quantiles corresponding to the data type and the neural network weights must be normalized to lie within this interval.

Given the impressive results achieved by {Guanaco} models, we examined whether there is any data leakage between the OASST1 dataset and the Vicuna benchmark prompts. Through fuzzy string matching and manual review of the closest matches, we found no overlapping prompts between these datasets.

    \textbf{Guanaco}: \texttt{Apologies, but as an AI language model, I do not have access to real-time data. My answers are generated from my existing knowledge and do not represent the current time. If you need to know the current time, I recommend checking a clock or using an online time source.}
\end{quote}

The Elo rankings presented in Table~\ref{tbl:elo_large} demonstrate that the {Guanaco} 33B and 65B models surpass all models except GPT-4 on the Vicuna and OA benchmarks, and their performance is on par with ChatGPT, consistent with the results in Table~\ref{tbl:vicuna_eval}. It is important to note that the Vicuna benchmark tends to favor open-source models, whereas the larger OA benchmark tends to favor ChatGPT. Additionally, Tables \ref{tbl:mmlu-llama} and \ref{tbl:vicuna_eval} reveal that the appropriateness of the finetuning dataset plays a crucial role in determining performance. For instance, finetuning Llama models on FLAN v2 yields particularly strong results on MMLU but leads to the lowest performance on the Vicuna benchmark (a similar pattern is observed with other models). This suggests a degree of orthogonality in current evaluation benchmarks: excelling on MMLU does not necessarily correlate with strong chatbot performance (as assessed by Vicuna or OA benchmarks), and vice versa.

In a similar vein, we observe that the default hyperparameters for fully finetuned baselines are not optimally tuned. We conduct a hyperparameter sweep over learning rates ranging from 1e-6 to 5e-5 and batch sizes between 8 and 128 to establish robust baselines. The results for finetuning the 7B LLaMA model on Alpaca are presented in Figure~\ref{fig:hyper}.

highlighting the necessity for further research into more reliable instruction adherence methods.

\begin{quote}
    \textbf{User}: What time is it?

\section{Background}
\paragraph{Block-wise k-bit Quantization} Quantization refers to the process of mapping an input representation containing more information to one with reduced information. Typically, this involves converting data from a type with a higher bit-width to one with fewer bits, such as changing 32-bit floating point numbers to 8-bit integers. To maximize usage of the available range in the low-bit data type, the input is generally rescaled into the target type’s range by normalizing with the absolute maximum value of the input elements, which are often organized as a tensor. For instance, quantizing a 32-bit floating point (FP32) tensor to an Int8 tensor with the range $[-127, 127]$ is performed as follows:

     \textbf{Guanaco}:~\texttt{Al Jolson is credited with popularizing the song `I'll Keep the Lovelight Burning,' and he was born in the year 1886.}
 \end{quote}

where doubleDequant$(\cdot)$ is given by:

\paragraph{k-bit \textsc{QLoRA} achieves parity with 16-bit full finetuning and 16-bit LoRA}

Next, we elaborate on the components of \textsc{QLoRA} and provide a formal definition of the method.

Overall, we believe \textsc{QLoRA} will have a predominantly positive influence by making the finetuning of high-quality LLMs far more accessible and straightforward for a broad audience.

    \textbf{Guanaco}: \texttt{James is aware that Abby put the pen in the desk drawer. However, James also knows he has relocated the pen to the bag, and he assumes Abby is unaware of this. Therefore, James might expect Abby to search for the pen in the desk drawer, as that is where she last left it. Nonetheless, this assumption could be mistaken if Abby saw James moving the pen.}
\end{quote}

    \textbf{Guanaco}: \texttt{Certainly, I’d be glad to assist! 1833 is a prime number, meaning it has exactly two factors: 1 and 1833. \\ \\ Thus, the prime factorization of 1833 is: \\ \\ 2\^{}1 * 3\^{}2 * 17\^{}1 \\ \\ I hope this is helpful! Feel free to ask if you have further questions.}
\end{quote}

We move forward to explore instruction tuning at scales that cannot be feasibly examined with full 16-bit finetuning on standard academic research hardware.

\section{QLoRA Compared to Standard Finetuning}
\label{sec:comp_to_std}

\begin{equation}
    \text{doubleDequant}(c_1^{\text{FP32}}, c_2^{\text{k-bit}}, \mathbf{W}^{\text{k-bit}}) = \text{dequant}(\text{dequant}(c_1^{\text{FP32}}, c_2^{\text{k-bit}}), \mathbf{W}^{\text{4bit}}) = \mathbf{W}^{\text{BF16}},
\end{equation}

\begin{equation}
    \mathbf{Y} = \mathbf{X}\mathbf{W} + s\mathbf{X}\mathbf{L}_1\mathbf{L}_2,
\end{equation}

We have explained the mechanism of QLoRA and its ability to greatly reduce memory requirements for model finetuning. The primary question now is whether QLoRA can achieve performance comparable to full-model finetuning. Additionally, we seek to examine the elements of QLoRA, including the effect of NormalFloat4 versus conventional Float4. The upcoming sections describe the experiments conducted to address these inquiries.

\paragraph{Experimental setup.} We evaluate three types of architectures (encoder, encoder-decoder, and decoder-only) and compare QLoRA against 16-bit adapter finetuning and full finetuning for models with sizes up to 3B parameters. Our tests cover GLUE \citep{wang2018glue} with RoBERTa-large \citep{liu2019roberta}, Super-NaturalInstructions (TKInstruct) \citep{wang2022super} using T5 \citep{t5}, and 5-shot MMLU \citep{hendrycksmeasuring} following finetuning LLaMA on Flan v2 \citep{longpre2023flan} and Alpaca \citep{alpaca}. To further explore the benefits of NF4 over other 4-bit formats, we adopt the experimental design of \citet{dettmers2022case} and assess post-quantization zero-shot accuracy and perplexity on various models (OPT~\citep{zhang2022opt}, LLaMA~\citep{touvron2023llama}, BLOOM~\citep{scao2022bloom}, Pythia~\citep{biderman2023pythia}) ranging from 125 million to 13 billion parameters.

Although paged optimizers are essential for tuning 33B/65B \textsc{QLoRA} models on a single 24/48GB GPU, we do not present exact measurements for Paged Optimizers because paging only happens when processing mini-batches with very long sequences, which is infrequent. Nonetheless, we analyze the runtime of paged optimizers on 65B models running on 48GB GPUs and observe that with a batch size of 16, paged optimizers achieve training speeds comparable to standard optimizers. Future investigations should aim to quantify and characterize conditions under which paging leads to slowdowns.

Nevertheless, {Guanaco} can falter even with straightforward problems if it does not decompose them into step-by-step reasoning, an issue recognized in prior work \cite{weichain}. For example, consider this exchange:

Another limitation is that we did not explore different bit-precision configurations, such as 3-bit base models, or alternative adapter techniques. Beyond LoRA, there exists a broad range of Parameter Efficient FineTuning (PEFT) methods that have demonstrated strong performance. However, it is unclear how well these methods scale to very large models. We selected LoRA due to its established robustness, but other adapters might achieve superior results. Since finetuning post-quantization appears to restore most of the information lost during quantization, this could potentially enable much more aggressive quantization approaches. For example, combining 3-bit GPTQ quantization of the base model with LoRA finetuning might reach performance levels comparable to full 16-bit finetuning after the finetuning process.

it produces an incorrect popularizer and an incorrect birth year (although the birth year provided is correct for the individual mentioned, Al Jolson).

\paragraph{Instruction Finetuning} To enable a pretrained LLM to better follow instructions within a prompt, instruction finetuning involves training the model using input-output pairs from diverse data sources, fine-tuning it to produce the desired output given the input prompt. Relevant approaches and datasets include MetaICL~\citep{min2021metaicl}, MetaTuning~\citep{zhong2021adapting}, InstructGPT~\citep{ouyang2022training}, FLAN~\citep{wei2021finetuned,chung2022scaling}, PromptSource~\citep{bach2022promptsource}, Super-NaturalInstructions~\citep{wang2022super,sanh2021multitask}, Self-instruct~\citep{wang2022self}, UnnaturalInstructions~\citep{honovich2022unnatural}, OPT-IML~\citep{iyer2022opt}, UnifiedSKG\citep{xie2022unifiedskg}, OIG/Chip2~\citep{laion2023}, Alpaca~\citep{alpaca}, Vicuna~\citep{vicuna2023}, Koala~\citep{koala_blogpost_2023}, and Self-instruct-GPT-4~\citep{peng2023instruction}.

Aligning with prior research on quantization \citep{dettmers2022case}, our results on MMLU and Elo suggest that, within fixed finetuning and inference resource constraints, it is advantageous to increase the size of the base model while lowering parameter precision. This underscores the efficiency gains enabled by \textsc{QLoRA}. Since our 4-bit finetuning experiments did not reveal any performance drop relative to full finetuning, it prompts further investigation into the precise nature of the performance-precision trade-off in QLoRA tuning, which we leave for future research.

Nonetheless, these deductions are often unreliable, with the model sometimes offering explanations based on assumptions that are illogical within the context, for example:

\paragraph{Refusal} Likewise, {Guanaco} occasionally declines to comply with instructions for apparently arbitrary reasons:

\paragraph{Data \& Training} It is important to highlight that the OASST1 dataset, which serves as the training data for {Guanaco} models, is multilingual. Additionally, the OA benchmark includes prompts in various languages. Future investigations could explore how multilingual training impacts the model's ability to follow instructions in languages other than English, and whether this factor accounts for the notable performance gap between the Vicuna-13B model (trained exclusively on English data) and the {Guanaco} 33B and 65B models on the OA benchmark.

To encourage further research, we release all model outputs along with annotations from both humans and GPT-4. We also open-source our codebase and CUDA kernels, integrating our methods into the Hugging Face transformers framework \citep{wolf2019huggingface}, thus ensuring broad accessibility. We provide a suite of adapters for models of sizes 7B, 13B, 33B, and 65B, trained on eight distinct instruction-following datasets, amounting to a total of 32 open-source finetuned models.

However, as questions become more obscure, {Guanaco} tends to become less reliable but remains confident. For example, in response to the following prompt from HotPotQA \citep{yang2018hotpotqa}:

\paragraph{Theory of Mind} {Guanaco} demonstrates notably strong Theory of Mind abilities \cite{nematzadeh2018evaluating, sap2022neural}. For example, the model provides a detailed and accurate response to the following query:

{Guanaco} is the only leading model in our evaluation that is not trained on proprietary data, given that the OASST1 dataset collection guidelines explicitly prohibit the use of GPT models. The next best model trained exclusively on open-source data is the Anthropic HH-RLHF model, which scores 30 percentage points lower than {Guanaco} on the Vicuna benchmark (refer to Table~\ref{tbl:vicuna_eval}). Overall, these findings illustrate that 4-bit \textsc{QLoRA} is an effective approach capable of producing state-of-the-art chatbots that rival ChatGPT. Moreover, our 33B {Guanaco} model can be trained on 24 GB consumer GPUs in under 12 hours. This paves the way for future work involving \textsc{QLoRA} tuning on specialized open-source datasets, enabling the creation of models that can compete with the best commercial models currently available.

\paragraph{4-bit NormalFloat outperforms 4-bit Floating Point}

\end{document}

\paragraph{Baselines} We benchmark our models against both research-oriented (Vicuna~\citep{vicuna2023} and Open Assistant~\citep{kopf2023openassistant}) and commercial (GPT-4 \citep{openai2023gpt}, GPT-3.5-turbo, and Bard) chatbot systems. The Open Assistant model is a LLaMA 33B variant finetuned via Reinforcement Learning from Human Feedback (RLHF) on the same OASST1 dataset used in our experiments. Vicuna represents a full finetuning of LLaMA 13B on proprietary user-shared conversations sourced from ShareGPT and thus constitutes a distillation derived from OpenAI GPT models.